\chapter{Supporting Conversion Experiments}
\label{app:conversion}

This appendix preserves the converter development and label-integrity checks that support the final procedure of Chapter~\ref{ch:data}. Historical versions are not the labels used for the main results.

\section{Development of the converter}
\label{sec:conversion-history}

The first converter (v1) chose the Anny bones by name. It gave \DataVOneBarMean\,mm on \texttt{downtown\_bar\_00};
on four test scenes the error was \DataVOneFloorMin{} to \DataVOneFloorMax\,mm on every frame, and all four actors
came out with the same body. A residual independent of the pose and a shape independent of the person mean
that the shape absorbed a fixed mismatch between the skeletons. Placing both models in their rest poses showed two
wrong choices that no name reveals. Anny's rig numbers the spine top-down, SMPL bottom-up, so SMPL's lumbar
\texttt{spine1} had been mapped to a bone \DataSpineMapError\,cm above it; and SMPL's hips had been mapped to
\texttt{pelvis.L/R}, \DataHipMapError\,cm from the hip joint, instead of \texttt{upperleg01.L/R}. Version 2 picks the
spine bones by height along the body axis and the hip bones by distance from the pelvis (Appendix~\ref{app:joints}).
Limbs and head fell to \DataVTwoLimbMin{} to \DataVTwoLimbMax\,mm, but the torso stayed at \DataVTwoTorsoMin{} to
\DataVTwoTorsoMax\,mm (Table~\ref{tab:fit-versions}).

The reason is that the rigs define their torso joints at different anatomical points: SMPL's pelvis joint sits
\DataSmplPelvisAboveHips\,cm above its hips, Anny's root at hip height and \DataAnnyRootFront\,cm in front. For
\DataVirtualJoints{} torso joints I therefore measure once, in the rest pose, the offset $\boldsymbol{\delta}_b$ from
the Anny bone head to the SMPL joint in the bone's own frame, and the fit moves the \emph{virtual joint}
$\mathbf{p}_b = \mathbf{h}_b + \mathbf{R}_b\boldsymbol{\delta}_b$ onto the SMPL joint, with $\mathbf{h}_b$ and
$\mathbf{R}_b$ the posed bone head and frame. The offsets are \DataOffsetPelvis\,cm (pelvis), \DataOffsetHip\,cm
(hips), \DataOffsetSpineMin{} to \DataOffsetSpineMax\,cm (spine) and \DataOffsetCollar\,cm (collars). With them
(v2.1) the torso fell to \DataVTwoOneTorsoMin{} to \DataVTwoOneTorsoMax\,mm and four varied test scenes to
\DataScenesMin{} to \DataScenesMax\,mm (one occluded frame \DataScenesWorstFrame\,mm). Age, however, came out at exactly \DataFreeAge{} for every actor in every
version: joints cannot tell one adult age from another. Version 2.2 fixes age at \DataFixedAge; on all
\DataThreeDPWPersonSeq{} person-sequences it fitted \DataOldFrames{} frames with a median error of
\DataOldFitMedian\,mm. By every measure the converter prints, these fits looked clean. They still held two of the
three defects of the next section.

\begin{table}[htbp]
\centering\small
% Draft main.tex:861-875 (tab:fit_versions), measured; jobs named in smpl_to_anny.py docstring. Cells kept literal.
% Correction: the draft's "fixed gender" for v2 hid that the fixed value was reversed until v2.3.
\begin{tabular}{@{}>{\raggedright\arraybackslash}p{5.2cm}rr>{\raggedright\arraybackslash}p{3.7cm}@{}}
\toprule
converter & mean (mm) & max & worst joints (mm) \\
\midrule
v1: correspondence by name, neutral SMPL targets & \DataVOneFencing & 37 & all joints, common floor \\
v2: correspondence by geometry, 3DPW joints, gender from metadata (reversed until v2.3) & \DataVTwoFencing & 30 & spine1 73, collars 55--59, spine3 58, pelvis 53 \\
v2.1: + virtual torso joints, bone mask, warm start & \DataVTwoOneFencing & 18.8 & shoulder 36, wrists 25--27 \\
\bottomrule
\end{tabular}
\caption[Converter versions on the fencing scene]{Mean joint error over all 24 joints on
\texttt{outdoors\_fencing\_01} (\DataFencingFrames{} frames, including full lunges) for three converter versions.
The geometric correspondence is worth about \DataGainCorrespondence\,mm, the virtual joints a further
\DataGainVirtual\,mm.}
\label{tab:fit-versions}
\end{table}

% -----------------------------------------------------------------------------
\section{Three label defects and the v2.3 re-fit}
\label{sec:data-defects}

The converter grades its own fit, so it cannot see errors in what it is asked to fit or in what happens to the fits
afterwards. I therefore wrote an independent check that reads the finished label files. On sampled images from every
sequence it compares the labels with 3DPW's joints, surface and 2D keypoints, confirms that the loader passes them on
unchanged, and tests gender, a label for every valid actor and the head turn. On the first label set (``old labels'',
v2.2 fits, about \DataOldCheckImages{} images) its grade was FAIL, with three defects (Table~\ref{tab:label-check}).

\paragraph{Defect 1: gender reversed.} The converter wrote 1.0 for men and 0.0 for women, MakeHuman's convention and
the reverse of Anny's, so every 3DPW man was labelled with a female body. Refitting with the flipped value lowered the
joint error for \DataGenderFlipMen{} of \DataGenderFlipMenOf{} men (median \DataGenderFlipMenBefore{} to
\DataGenderFlipMenAfter\,mm); for women it raised the joint error (\DataGenderFlipWomenBefore{} to
\DataGenderFlipWomenAfter\,mm) but brought the surface closer for \DataGenderFlipWomenSurface{} of
\DataGenderFlipWomenOf.

\paragraph{Defect 2: missing actor.} In \DataMissingActor{} of \DataTwoPersonChecked{} sampled two-person images,
one valid actor had no label. Each person's every-third-frame grid started at that person's own first valid frame, so
after an invalid gap the two actors sat on different phases. The detection loss then treats a visible, real person as
background.

\paragraph{Defect 3: missing labels.} The training split had \DataOldTrainImages{} label images where the fits implied
\DataRebuiltTrainImages{} (\DataOldTrainMissingPct\% missing). The image files of \DataMissingImageSeq{} sequences
were incomplete on the cluster (\texttt{courtyard\_captureSelfies\_00}: \DataSelfiesFilesBefore{} of
\DataSelfiesFilesAfter{} files), and the builder skipped frames without an image file without counting them. A
rebuild after re-extracting the images recovered every label (Table~\ref{tab:label-sets}, ``rebuilt''), but still
had defects 1 and 2, so I never trained on it.

\paragraph{The re-fit.} Version 2.3 adds the two switches named above; their defaults reproduce v2.2 exactly. I
re-fitted all \DataThreeDPWPersonSeq{} people with both switches on, and the post-check passed on all \DataRefitSeq{}
sequences (Table~\ref{tab:label-check}, middle block). \DataRefitWorseSeq{} single-person outdoor sequences became
\DataRefitWorseMin{} to \DataRefitWorseMax\,mm worse, likely from a different warm start, which I did not test.

\paragraph{The v23 labels.} I call the labels built from the v2.3 fits the v23 labels: \DataThreeDPWTotalImages{}
images and \DataThreeDPWTotalPeople{} people (Table~\ref{tab:label-sets}). They have fewer images than the rebuilt set
but the same people, because both actors now share each frame. No person was dropped for fit error or box size.
\DataCrossCountryFramesNew{} frames were dropped for visibility, all from \texttt{outdoors\_crosscountry\_00}: with
3DPW's own camera, its ground-truth joints sit at a depth of $\DataCrossCountryDepth$\,m, behind the camera. The 3D
fit is fine; image labels for this sequence cannot exist with this camera, and no label set ever contained it.



\paragraph{The check on the v23 labels.} The same check on the v23 labels (\DataCheckPeople{} people in
\DataCheckImages{} images, \DataCheckSeq{} sequences) moved the grade from FAIL to WARN
(Table~\ref{tab:label-check}). The median body-joint error fell from \DataCheckJointsOld{} to
\DataCheckJointsNew\,mm, flipping the gender now makes the fit worse, and no sampled two-person image misses an actor.
The head turn relative to the chest is now kept; in the old test labels it was mostly lost
(\DataCheckHeadTestLabelOld$^\circ$ against \DataCheckHeadTestSmplOld$^\circ$ in SMPL). Two rows stay at WARN, the
limb-length ratios and the people with a joint more than \DataCheckFarMm\,mm off, mostly at the hip, likely because
Anny's and SMPL's limb proportions and hip joints differ (Table~\ref{tab:label-check}).





% -----------------------------------------------------------------------------
